%% ma: L12-B19-0031
%% lop: 12
%% chuong: 6
%% bai: 19
%% dang: 12.19.2
%% loai: TL
%% muc: VD
%% dap_an: "-"
%% nguon: "(NBV)-10 ĐỀ KIỂM TRA HK2-TOÁN 12-DUNG CHUNG 3 SÁCH.pdf (H02, Đề số 2), Phần 4 Câu 3 (tự luận)"
%% kien_thuc: "Bài 19 (công thức xác suất toàn phần)"
%% trang_thai: chinh-thuc
%% ghi_chu: "Câu tự luận. Đáp án gốc 103/160 khớp. Bỏ ảnh hai chuồng thỏ (ảnh trang trí)."
\begin{ex}
Có hai chuồng thỏ, chuồng thứ I có $5$ con thỏ đen và $10$ thỏ trắng, chuồng II có $3$ con thỏ trắng và $7$ con thỏ đen. Từ chuồng thứ II ta bắt ngẫu nhiên một con thỏ chuyển sang chuồng thứ I, rồi sau đó lại bắt ngẫu nhiên một con thỏ từ chuồng thứ I, tính xác suất để bắt được một con thỏ trắng.
\loigiai{Gọi $B$: ``con thỏ chuyển từ chuồng II sang là thỏ trắng'', $A$: ``con thỏ bắt từ chuồng I là thỏ trắng''. Khi đó $P(B)=\dfrac3{10}$, $P(\overline B)=\dfrac7{10}$.

Nếu con chuyển sang là thỏ trắng, chuồng I có $16$ con, trong đó $11$ thỏ trắng: $P(A\mid B)=\dfrac{11}{16}$. Nếu là thỏ đen, chuồng I có $16$ con, trong đó $10$ thỏ trắng: $P(A\mid\overline B)=\dfrac{10}{16}$.

Theo công thức xác suất toàn phần: $P(A)=\dfrac3{10}\cdot\dfrac{11}{16}+\dfrac7{10}\cdot\dfrac{10}{16}=\dfrac{103}{160}$.}
\end{ex}
